\documentclass{article}
\usepackage{graphicx}
\usepackage{listings}
\usepackage{color}

\definecolor{dkgreen}{rgb}{0,0.6,0}
\definecolor{gray}{rgb}{0.5,0.5,0.5}
\definecolor{mauve}{rgb}{0.58,0,0.82}

\lstset{frame=tb,
  language=Python,
  aboveskip=3mm,
  belowskip=3mm,
  showstringspaces=false,
  columns=flexible,
  basicstyle={\small\ttfamily},
  numbers=none,
  numberstyle=\tiny\color{gray},
  keywordstyle=\color{blue},
  commentstyle=\color{dkgreen},
  stringstyle=\color{mauve},
  breaklines=true,
  breakatwhitespace=true,
  tabsize=3
}

\title{Practical Work 3: MPI File Transfer}
\author{Vu Nhat Anh - 22BA13036}
\date{\today}

\begin{document}

\maketitle

\section{Choice of MPI Implementation}
This implementation uses \textbf{mpi4py}, a Python binding for the MPI standard. The choice was made for the following reasons:
\begin{itemize}
    \item It provides a clean, Pythonic interface to standard MPI operations.
    \item It supports automatic serialization of Python objects, simplifying data transfer.
    \item The concepts and APIs align with MPI implementations in C/C++, making the knowledge transferable.
\end{itemize}

\section{MPI Design}
MPI follows the SPMD (Single Program, Multiple Data) model, where the same program runs on multiple processes, each identified by a unique rank. The behavior of each process is determined by its rank value.

For file transfer, we use point-to-point communication:
\begin{itemize}
    \item \textbf{Rank 0}: Acts as the sender. Reads the file and transmits it to Rank 1.
    \item \textbf{Rank 1}: Acts as the receiver. Waits for the file data and saves it locally.
\end{itemize}

Messages are distinguished using tags: tag 0 for the filename, tag 1 for the file content.

\section{System Organization}
The system requires exactly two MPI processes. Both processes execute the same code, but their behavior differs based on rank:
\begin{itemize}
    \item Process 0 initiates the transfer by calling \texttt{comm.send()}.
    \item Process 1 waits for incoming data using \texttt{comm.recv()}.
\end{itemize}

The program is executed using: \texttt{mpiexec -n 2 python transfer.py <filename>}

\section{Implementation}
The implementation uses \texttt{MPI.COMM\_WORLD} for communication between processes:
\begin{lstlisting}
comm = MPI.COMM_WORLD
rank = comm.Get_rank()

if rank == 0:
    # Sender
    with open(filename, 'rb') as f:
        content = f.read()
    comm.send(filename, dest=1, tag=0)
    comm.send(content, dest=1, tag=1)
    
elif rank == 1:
    # Receiver
    filename = comm.recv(source=0, tag=0)
    content = comm.recv(source=0, tag=1)
    with open('received_' + filename, 'wb') as f:
        f.write(content)
\end{lstlisting}

\end{document}
